\begin{lemma}
Let $g\in\mathsf{IncSeq}$ be nonidentity.  There is a prefix-continuous map
$\rho:\mathsf{IncSeq}\to\mathsf{IncSeq}$ satisfying
\[
\rho(f\circ g)=\rho(f)\circ\mathsf{SuccSeq}
\]
for every $f$.  Equivalently, right composition by $g$ continuously
homomorphically maps to the ordinary shift.
\end{lemma}
\begin{proof}
By \noderef{FirstMovedPoint}, choose the least moved point $k_g$ with
$k_g<g(k_g)$.  Define $G(n)=g^{[n]}(k_g)$.  The node
\noderef{OrbitEmbedding} gives an increasing embedding $G$.  Define
$\rho(f)=f\circ G$.  Composition with the fixed finite-prefix map $G$ is
prefix-continuous: to determine the first $m$ values of $\rho(f)$ it is
enough to know the first $G(m-1)+1$ values of $f$.  It is an increasing
embedding because both $f$ and $G$ are.

For every $n$,
\[
\rho(f\circ g)(n)
 =f(g(G(n)))
 =f(G(n+1))
 =\rho(f)(n+1).
\]
The middle equality is the defining orbit identity.  Hence
$\rho(f\circ g)=\rho(f)\circ\mathsf{SuccSeq}$, which is the required
continuous homomorphism.
\end{proof}
